\begin{textsample}{POS Dim 4 – 2010\_c6 – Score -1.00 – 2010\_c6/enade\_2010\_tec\_agro\_87acf19614.txt}  \label{ex:f4_pos_312}
CAPÍTULO 3 ANÁLISE TÉCNICA DA PROVA Este capítulo tem por objetivo apresentar o desempenho dos estudantes de Tecnologia em Agronegócio no ENADE/2010.
Para isso, foram calculadas as estatísticas básicas da prova em seu todo, bem como as estatísticas dos componentes relacionadas à Formação Geral e ao de Conhecimento Específico da Área.
Nas tabelas, são apresentados o tamanho da população inscrita e de presentes, e as seguintes estatísticas das notas⁷ : média do desempenho na prova, erro padrão da média, desvio padrão, nota mínima, mediana, nota máxima, coeficiente de assimetria e de variação.
As estatísticas apresentadas neste capítulo contemplam o total de estudantes e, separadamente, os Ingressantes e os Concluintes.
Foram calculadas tendo -se em vista as seguintes agregações : ( a ) as Grandes Regiões e o país como um todo ; ( b ) a Categoria Administrativa ; e ( c ) a Organização Acadêmica.
Em relação aos gráficos de distribuição de notas, o intervalo considerado foi de 10 unidades, aberto à esquerda e fechado à direita, com exceção do primeiro intervalo, [ 0 ; 10 ], fechado em ambos os extremos.
Para os gráficos de distribuição das notas das questões discursivas, foram consideradas mais duas categorias : questão em branco e nota zero. 3.1 ESTATÍSTICAS BÁSICAS DA PROVA 3.1.1 Estatísticas Básicas Gerais A Tabela 3.1 apresenta as estatísticas básicas da prova por grupo de estudantes.
A população total de estudantes regularmente inscritos⁸ para a prova foi 2.311.
Destes, 1.795 estiveram presentes, sendo de 22,3 \% o índice de não comparecimento.
A abstenção foi maior entre Ingressantes ( 27,9 \% ) do que entre Concluintes ( 9,3 \% ).
A nota média de toda a prova foi 39,6, sendo que os Ingressantes obtiveram média mais baixa, 38,1, que os Concluintes, cuja média foi 42,3.
O desvio padrão foi 13,4, sendo o do grupo de Ingressantes um pouco menor ( 13,1 ) do que o do grupo dos Concluintes ( 13,6 ) ; os coeficientes de variação ( 34,4 \% no grupo dos Ingressantes e 32,2 \% no dos Concluintes ) indicam menor dispersão relativa nas notas de Concluintes.
A nota máxima foi 81,4, obtida por Ingressantes, sendo superior a de Concluintes ( 78,0 ). ⁷ Algumas dessas estatísticas estão definidas no Capítulo 1. ⁸ As estatísticas das notas apresentadas consideram apenas os estudantes regularmente inscritos e com resultados validados por ausência de erro administrativo.
Uma comparação entre o desempenho de Ingressantes e Concluintes também pode ser realizada a partir dos dados apresentados no Gráfico 3.1, com a distribuição de notas finais.
Destaca -se que as proporções mais elevadas de notas, considerando -se o intervalo posterior à nota 40,0, foram atingidas principalmente pelos Concluintes, cuja concentração das notas está à direita dos Ingressantes, com a exceção do último grupo com notas : ( 80 a 90 ].
Já nas faixas de notas mais baixas, inverteu -se o quadro identificado nas faixas de notas mais altas.
Essa característica é corroborada pelo valor da nota mediana, 37,2 no grupo dos Ingressantes, abaixo da encontrada no grupo dos Concluintes, 41,9, transparecendo uma superioridade de desempenho dos Concluintes em relação aos Ingressantes na prova em tela ( Tabela 3.1 ).
Os coeficientes de assimetria indicam um comportamento simétrico ( 0,0 ) para os Concluintes, e positivo ( 0,2 ) para Ingressantes.
O Gráfico 3.2 apresenta informações referentes à nota final de Ingressantes e Concluintes, comparando -se os resultados de acordo com as Grandes Regiões do país, a Categoria Administrativa e a Organização Acadêmica.
Considerando -se as notas médias dos alunos em cada Grande Região, observa -se que, em relação às notas aferidas em âmbito nacional ( média de 38,1 para Ingressantes e de 42,3 para Concluintes ), existe diferença estatisticamente significativa ao nível de 95 \%, para Ingressantes, entre a menor média, a obtida na região Sudeste ( 35,7 ) e a maior, na região Sul ( 39,5 ), nesse grupo.
Entre os Concluintes, já não existe diferença estatisticamente significativa ao nível de 95 \% entre as médias.
A maior média de notas foi obtida pela região Norte ( 44,2 ), ao passo que a menor foi identificada na região Centro-Oeste ( 41,5 ).
Ressalta -se que não houve participação de Concluintes na região Nordeste.
As demais médias de notas foram 41,8 para a região Sudeste e 42,7 para a região Sul.
Levando -se em conta os grupos dos estudantes em Categorias Administrativas, observa -se que a pontuação média mais elevada entre os Concluintes foi encontrada em instituições públicas ( 46,0 ), com desempenho acima da média nacional.
Já os Concluintes provenientes de instituições privadas obtiveram média 41,2, abaixo da média nacional.
Quanto aos Ingressantes, as notas mais elevadas, também acima da média nacional para esse grupo de estudantes, ocorreram nas instituições públicas ( 42,5 ).
Os Ingressantes vinculados a instituições privadas obtiveram nota média 36,8, abaixo da média nacional.
Tendo como foco as notas médias dos estudantes agrupados por Organização Acadêmica, distribuídas entre Universidades, Centros Universitários e Faculdades, para os Concluintes vinculados a Universidades, a média de 44,6 ficou acima da média nacional.
Pode -se verificar a mesma realidade entre os Ingressantes de Universidades, cuja nota média foi 39,7, também acima da nacional para esse grupo.
Entretanto, Ingressantes de Centros Universitários \textbf{ultrapassaram} os de Universidade, com nota média 40,3.
Nesse mesmo tipo de organização acadêmica, os Concluintes obtiveram a mesma nota média dos estudantes de Universidades, ou seja, 44,6.
Quanto aos provenientes de Faculdades, Ingressantes alcançaram nota média 35,7, enquanto Concluintes atingiram 40,9, ambas abaixo da média nacional. 3.1.2 Estatísticas Básicas no Componente de Formação Geral A Tabela 3.2 apresenta as estatísticas básicas em relação ao componente da prova que avalia a Formação Geral dos estudantes.
Os Concluintes obtiveram desempenho médio de 38,0, que foi superior ao dos Ingressantes, 36,2.
Quanto à variabilidade das notas nos dois grupos, o desvio padrão das notas obtidas foi semelhante : 16,6 para Ingressantes e 16,2 para Concluintes.
Os dois grupos alcançaram notas máximas aproximadas, sendo 86,5 a maior nota obtida por pelo menos um Ingressante, e 85,0 a maior nota entre os Concluintes.
O Gráfico 3.3 propicia a comparação entre Concluintes e Ingressantes relativa ao desempenho no componente de Formação Geral.
Nesse componente, os Concluintes apresentaram desempenho superior.
Dos Ingressantes, um pouco mais de 20,0 \% obtiveram notas na faixa entre 20 e 30 pontos, sendo essa a faixa modal.
Nas faixas de notas seguintes, mais elevadas, verifica -se queda gradual de ocorrência de estudantes Ingressantes.
Entre os Concluintes, houve percentual maior de estudantes no intervalo entre 30 e 40 pontos, acima dos 20,0 \%.
Também foi observada queda gradual no desempenho, nos intervalos de notas mais elevados, com 7,0 \% desses estudantes tendo atingido a faixa de 60 a 70 e 1,9 \% a de notas entre 70 e 80 pontos.
Apenas 4,1 \% dos Concluintes tiveram desempenho na menor faixa de notas, que vai de zero a 10 pontos.
O coeficiente de assimetria foi positivo ( 0,2 ) para Ingressantes, e simétrico ( 0,0 ) para Concluintes.
No Gráfico 3.4 são apresentadas as informações referentes ao desempenho dos Ingressantes e dos Concluintes no componente de Formação Geral, em diferentes agregações como Grande Região do país, Categoria Administrativa e Organização Acadêmica.
Considerando -se as notas médias dos estudantes segundo cada Grande Região, observa -se que no grupo de Ingressantes existe diferença estatisticamente significativa ao nível de 95 \% entre as médias nas regiões Nordeste ( 42,7, a mais alta ) e Sudeste ( 32,0, a mais baixa ).
No grupo de Concluintes, a nota mais elevada ocorreu na região Norte, com média 41,0.
A diferença entre notas mais altas e mais baixas não foi estatisticamente significativa ao nível de 95 \% : a menor nota foi 36,8, correspondeu aos Concluintes da região Sudeste.
Há que se considerar, mais uma vez, que não houve Concluintes da região Nordeste participando do exame.
Destaca -se que a região Sudeste ficou abaixo da média nacional no componente de Formação Geral.
Tendo como foco as Categorias Administrativas, as instituições públicas obtiveram médias mais elevadas do que as privadas : para Concluintes, notas médias de 41,8 em cursos de IES públicas e de 36,9 nos das privadas.
Para os Ingressantes, as notas foram, respectivamente, 42,0 ( superior a de Concluintes ) e 34,6.
Em comparação à média nacional, as notas de Ingressantes e de Concluintes vinculados às instituições públicas foram superiores, ao passo que, nas instituições privadas, foram inferiores.
Considerando -se o tipo de Organização Acadêmica, não se nota estatisticamente significativa ao nível de 95 \% entre a maior e a menor média.
Nas Universidades, houve desempenho melhor do que em Centros Universitários e Faculdades.
Neste primeiro tipo de Organização Acadêmica, Ingressantes tiveram nota média 38,0, enquanto Concluintes obtiveram 40,6.
Nos Centros Universitários e Faculdades, os valores foram mais baixos : os Ingressantes tiveram média 37,4 nos Centros Universitários e 34,4 nas Faculdades ; entre Concluintes, a média foi 40,0 para Centros Universitários e 36,6 para as Faculdades. 3.1.3 Estatísticas Básicas do Componente de Conhecimento Específico A Tabela 3.3 apresenta as estatísticas básicas referentes à parte do componente de Conhecimento Específico da Área de Tecnologia em Agronegócio.
A média do desempenho foi de 40,4.
Os Concluintes apresentaram desempenho médio ( 43,7 ) superior ao dos Ingressantes ( 38,7 ) ; as medidas de dispersão – desvio padrão e amplitude geral das notas – são ligeiramente menores no grupo dos Ingressantes do que no grupo dos Concluintes.
Os coeficientes de variação, 36,2 \% no grupo dos Ingressantes e 33,9 \% no dos Concluintes, indicam dispersão relativa menor nas notas desses últimos.
As notas máximas obtidas por Concluintes e Ingressantes foram iguais a 84,8 e 81,0, respectivamente, notas um pouco inferiores àquelas encontradas no componente de Formação Geral ( 85,0 e 86,5, respectivamente ).
As medianas das distribuições das notas dos Ingressantes e dos Concluintes foram 38,7 e 43,2, respectivamente, enfatizando a superioridade dos Concluintes neste componente da prova.
Assim como os Gráficos 3.1 e 3.3, o Gráfico 3.5, apresentado a seguir, proporciona uma comparação do desempenho entre Concluintes e Ingressantes, em relação ao componente de Conhecimento Específico.
Dentre as três comparações feitas, observa -se que, embora a distribuição das notas seja semelhante à apresentada na prova como um todo, neste componente, a diferença de desempenho entre Concluintes e Ingressantes é mais acentuada, destacando -se o desempenho mais elevado dos primeiros.
Claramente a distribuição das notas dos Ingressantes, mais concentradas à esquerda em relação aos Concluintes, evidencia o percentual mais alto de notas nas faixas mais baixas para este grupo.
A maior concentração de Concluintes e Ingressantes está no intervalo compreendido entre ( 30 a 40 ] pontos, com respectivamente 24,8 e 29,3 \% dos estudantes do grupo.
Assim como os Gráficos 3.2 e 3.4, o Gráfico 3.6 apresenta uma comparação dos resultados em relação à Grande Região do país, à Categoria Administrativa e à Organização Acadêmica, agora levando em conta o desempenho dos Ingressantes e dos Concluintes no componente de Conhecimento Específico da prova.
Ao se observar as notas médias dos estudantes, segundo cada Grande Região do país, verifica -se que, entre os Concluintes, a mais elevada foi encontrada na região Norte ( 45,2 ), maior que a média nacional ( 43,7 ).
Em relação aos Ingressantes, a região Sul apresenta a média mais elevada ( 40,2 ), seguida da região Norte, com média igual a 38,7.
A média mais baixa entre os Concluintes ( 42,0 ) foi verificada na região Centro-Oeste, enquanto para Ingressantes ( 34,9 ), na Nordeste.
Quanto à Categoria Administrativa, observa -se comportamento semelhante àquele da parte relativa à Formação Geral e à prova como um todo, ou seja, as médias mais elevadas tanto entre os Concluintes quanto entre os Ingressantes foram encontradas em instituições de ensino públicas.
Nesse tipo de IES, a nota média de Concluintes foi 47,4, enquanto a de Ingressantes, 42,6.
Tais médias estão situadas acima da média nacional.
Nas instituições privadas, as médias foram inferiores em relação às públicas e às médias nacionais.
Concluintes obtiveram média 42,7 e Ingressantes, 37,6.
Em relação à Organização Acadêmica, o desempenho dos Concluintes foi superior ao dos Ingressantes nas três categorias.
Nos Centros Universitários, tais estudantes obtiveram a média mais elevada, igual a 46,2.
Um pouco abaixo estiveram as Universidades com média 45,9.
Os Concluintes de Faculdades obtiveram média 42,4.
Os Ingressantes de Centros Universitários também tiveram nota um pouco mais elevada ( 41,3 ), seguidos pelos estudantes de Universidades ( 40,3 ) e, finalmente, os de Faculdades ( 36,1 ). 3.2 ANÁLISE DAS QUESTÕES OBJETIVAS 3.2.1 Componente de Formação Geral A Tabela 3.4 apresenta as estatísticas básicas relativas às oito questões objetivas do componente da prova que abrange a Formação Geral dos estudantes.
Como pode ser observado, os estudantes Concluintes obtiveram desempenho superior, com nota média de 41,0, em relação aos alunos Ingressantes, cuja nota média foi 38,5.
As medidas de dispersão nos dois grupos de alunos foram similares, como indicam os desvios padrões, assim como a amplitude total das notas, sendo que a nota mínima foi zero e a máxima, 100,0, para os dois grupos de alunos avaliados.
Os coeficientes de variação, 48,3 \% no grupo dos Ingressantes e 45,1 \% no dos Concluintes, indicam menor dispersão relativa nas notas desses últimos.
A mediana foi igual a 37,5 nos dois grupos.
O coeficiente de assimetria do total de estudantes revelou assimetria positiva ( 0,1 ).
A Tabela 3.5 apresenta as classificações das questões objetivas do componente de Formação Geral, segundo o índice de facilidade.
As questões objetivas da prova foram assim avaliadas : das oito questões, nenhuma teve o índice de facilidade classificado como muito fácil.
Uma questão foi tida como fácil, por ter índice de acertos situado na faixa entre 0,61 e 0,85 ( de 61,0 \% a 85,0 \% de acertos ).
Três questões foram consideradas de facilidade média, situando -se no intervalo entre 0,41 e 0,60 do índice de facilidade, ou seja, houve entre 41,0 \% e 60,0 \% de acertos.
Duas questões foram classificadas como difíceis, situando -se no intervalo entre 0,16 e 0,40.
Por fim, duas questões apresentaram menos de 0,15 de acertos, razão pela qual foram classificadas como muito difíceis.
A Tabela 3.6 descreve os resultados das análises das questões objetivas relativas à Formação Geral, segundo o poder de discriminação, utilizando -se para tal, do índice de discriminação ponto bisserial.
Cinco das oito questões apresentaram índices acima de 0,40 e, assim, foram classificadas como muito boas nesta característica para esse grupo de alunos ; uma questão teve bom índice de discriminação, entre 0,30 e 0,39.
Foram classificadas como de média discriminação duas das questões.
Nenhuma delas teve índice de discriminação fraco.
Na sequência, na Tabela 3.7, estão especificadas, por questão, os resultados da análise relativa aos índices de facilidade e de discriminação de cada questão no componente relativo à Formação Geral da prova da Tecnologia em Agronegócio.
O índice de facilidade variou de 0,11 a 0,69, e o de discriminação, de 0,22 a 0,56.
As questões com índices de discriminação muito bom, de números 1, 3, 4, 5 e 8, figuraram entre as com maior poder de discriminação.
A questão 8 foi a que apresentou maior poder discriminatório, com índice de 0,56, e foi também uma das mais fáceis, com uma proporção de 60,0 \% de acertos.
A questão de número 6 apresentou índice de discriminação de 0,22, e um índice de facilidade de 0,11, ou seja, um quantitativo mais reduzido de estudantes conseguiu resolvê -la, dentro do universo de participantes.
O mesmo foi verificado para a questão 7.
A Tabela 3.6 descreve os resultados das análises das questões objetivas relativas à Formação Geral, segundo o poder de discriminação, utilizando -se para tal, do índice de discriminação ponto bisserial.
Cinco das oito questões apresentaram índices acima de 0,40 e, assim, foram classificadas como muito boas nesta característica para esse grupo de alunos ; uma questão teve bom índice de discriminação, entre 0,30 e 0,39.
Foram classificadas como de média discriminação duas das questões.
Nenhuma delas teve índice de discriminação fraco.
Na sequência, na Tabela 3.7, estão especificadas, por questão, os resultados da análise relativa aos índices de facilidade e de discriminação de cada questão no componente relativo à Formação Geral da prova da Tecnologia em Agronegócio.
O índice de facilidade variou de 0,11 a 0,69, e o de discriminação, de 0,22 a 0,56.
As questões com índices de discriminação muito bom, de números 1, 3, 4, 5 e 8, figuraram entre as com maior poder de discriminação.
A questão 8 foi a que apresentou maior poder discriminatório, com índice de 0,56, e foi também uma das mais fáceis, com uma proporção de 60,0 \% de acertos.
A questão de número 6 apresentou índice de discriminação de 0,22, e um índice de facilidade de 0,11, ou seja, um quantitativo mais reduzido de estudantes conseguiu resolvê -la, dentro do universo de participantes.
O mesmo foi verificado para a questão 7.
O Gráfico 3.7 analisa o comportamento da questão de número 8 de Formação Geral.
Trata -se uma das questões mais fáceis e a com o maior índice de discriminação desse componente da prova.
No gráfico a seguir, cada uma das cinco curvas representa o percentual de respostas em determinada alternativa da questão, em função da nota dos estudantes nesta parte da prova ( Formação Geral/Múltipla Escolha ) antes de possíveis eliminações pelo critério do ponto bisserial.
A curva em rosa corresponde à alternativa B, a correta para este item.
Assim, observa -se que entre os estudantes com notas mais baixas, desempenho mais fraco neste componente do exame, a situação mais frequente foi a escolha de uma das alternativas incorretas : a alternativa C ( em verde ) ou E ( em roxo ).
Na medida em que a nota aumenta, indicando desempenho melhor nesta parte da prova, aumenta, concomitantemente, a proporção de estudantes que selecionaram a alternativa correta B, atingindo 100 \% para as notas acima de 80,0.
Essa análise permite verificar como a questão discriminou entre os grupos de desempenho, justificando o alto índice obtido na questão.
Os gráficos relativos às demais questões de Formação Geral constam do Anexo I. 3.2.2 Componente de Conhecimento Específico A Tabela 3.8 apresenta as estatísticas básicas em relação às questões objetivas do componente de Conhecimento Específico da prova.
Nesse conjunto de questões, os estudantes Concluintes ( 48,3 ) obtiveram nota média superior aos Ingressantes ( 42,9 ).
Quanto à variabilidade das notas, o desvio padrão no grupo dos Ingressantes foi 15,3 e no grupo dos Concluintes, 16,1.
Os coeficientes de variação foram 35,7 \% e 33,3 \%, respectivamente, indicando menor dispersão relativa nas notas dos Concluintes.
Com relação à mediana, no grupo dos Ingressantes, esta foi igual a 40,9, estando pouco abaixo da média desse grupo, enquanto no grupo dos Concluintes, a mediana foi igual a 50,0, superando um pouco a média desse grupo.
Os coeficientes de assimetria foram 0,1 para Ingressantes e simétrico ( 0,0 ) para Concluintes.
A Tabela 3.9 apresenta as classificações das questões objetivas da parte da prova relativa ao componente de Conhecimento Específico, segundo o índice de facilidade.
A partir dos índices obtidos, pode -se concluir que, a prova teve um índice de facilidade mediano.
Menos da metade das questões objetivas da prova foram tidas como pelo menos difíceis : das 27 questões, onze foram classificadas como difíceis e duas como muito difíceis.
Não houve questões classificadas como muito fáceis, ao passo que três foram tidas como fáceis, na faixa de 0,61 a 0,85 do índice de facilidade, e onze consideradas médias, entre 0,41 e 0,60.
Na Tabela 3.10 são apresentados os índices de discriminação das questões objetivas do componente de Conhecimento Específico da prova.
Observa -se, segundo os índices estabelecidos, que a capacidade de discriminação das questões objetivas desse componente foi classificada, para nove das 27 questões, como boa, enquanto três delas tiveram índice de discriminação muito bom.
Observa -se, portanto, que 12 das 27 questões proporcionaram índices de discriminação bons ou muito bons.
Dez delas foram A Tabela 3.9 apresenta as classificações das questões objetivas da parte da prova relativa ao componente de Conhecimento Específico, segundo o índice de facilidade.
A partir dos índices obtidos, pode -se concluir que, a prova teve um índice de facilidade mediano.
Menos da metade das questões objetivas da prova foram tidas como pelo menos difíceis : das 27 questões, onze foram classificadas como difíceis e duas como muito difíceis.
Não houve questões classificadas como muito fáceis, ao passo que três foram tidas como fáceis, na faixa de 0,61 a 0,85 do índice de facilidade, e onze consideradas médias, entre 0,41 e 0,60.
Na Tabela 3.10 são apresentados os índices de discriminação das questões objetivas do componente de Conhecimento Específico da prova.
Observa -se, segundo os índices estabelecidos, que a capacidade de discriminação das questões objetivas desse componente foi classificada, para nove das 27 questões, como boa, enquanto três delas tiveram índice de discriminação muito bom.
Observa -se, portanto, que 12 das 27 questões proporcionaram índices de discriminação bons ou muito bons.
Dez delas foram classificadas como médias e cinco como fracas, sendo quinze, por conseguinte, a quantidade de questões nos dois patamares mais baixos de discriminação.
Constata -se, assim, que a prova – no que se refere ao componente de Conhecimento Específico – possuía capacidade mediana de discriminar entre aqueles que dominam ou não o conteúdo.
As questões classificadas com índice de discriminação fraco ( questões 17, 18, 29, 32 e 36 ) foram eliminadas do cômputo das notas.
Na Tabela 3.11 estão especificados, por questão, os resultados da análise relativa aos índices de facilidade e de discriminação de cada uma das questões do componente de Conhecimento Específico da prova da Tecnologia em Agronegócio.
Dentre as questões que alcançaram os maiores índices de discriminação, muito bom, as de números 15, 24 e 30 foram classificadas como de dificuldade média.
Em termos de proporção, tais questões variaram de 0,46 a 0,49 quanto aos acertos.
Seus índices de discriminação situaram -se de 0,41 a 0,46.
A questão de número 18 foi a mais difícil dentre as 27 questões específicas, com baixo índice de facilidade, de 0,13, o que, em termos percentuais, corresponde a apenas 13,0 \% de acertos.
Essa questão apresentou poder de discriminação igualmente baixo, 0,14, o que comprova ter sido esta uma das mais difíceis para os estudantes.
Destaca -se, também, a questão de número 32, que teve o índice de facilidade de 0,14 o que, em termos percentuais, corresponde a 14,0 \% de estudantes que responderam acertadamente, obtendo também índice de discriminação fraco, 0,14.
O Gráfico 3.8 analisa o comportamento da questão 30 do componente de Conhecimento Específico.
Esta foi uma das questões mais fáceis nesta parte da prova, apresentando índice de facilidade 0,47, ou seja, 47,0 \% dos estudantes assinalaram acertadamente a opção A, correspondente ao gabarito.
Seu índice de discriminação foi igual a 0,46, classificado como muito bom.
Neste gráfico, cada uma das cinco curvas representa o percentual de respostas em determinada alternativa da questão 30, em função da nota dos estudantes nesta parte da prova antes de possíveis eliminações de questões pelo critério do ponto bisserial.
A alternativa correta A, representada no gráfico pela curva em azul, foi escolhida em maiores proporções pelos alunos com desempenho melhor nesse componente do exame, enquanto as alternativas incorretas, também denominadas distratores, foram selecionadas principalmente por aqueles com notas mais baixas, ainda que aqueles com nota zero, na sua quase totalidade tenham deixado em branco ou marcaram mais de uma alternativa, comportamento considerado inválido.
A proporção de alunos que selecionaram a resposta correta A aumenta gradativamente, com a nota chegando a atingir 100 \% para médias acima de 80, enquanto a proporção dos que escolhem alternativas incorretas decai a partir da primeira nota não nula, como função da nota nesta parte da prova.
Os gráficos relativos às demais questões do componente de Conhecimento Específico constam do Anexo I. 3.3 ANÁLISE DAS QUESTÕES DISCURSIVAS 3.3.1 Componente de Formação Geral A análise dos resultados de desempenho dos estudantes de Tecnologia em Agronegócio, nas duas questões discursivas relativas à Formação Geral, encontra -se na Tabela 3.12 e no Gráfico 3.9.
Observa -se, a partir destes, que as notas médias foram mais baixas nesse conjunto do que no conjunto das objetivas.
A média geral dos Ingressantes nas questões objetivas de Formação Geral foi 38,5 ( Tabela 3.4 ), ao passo que nas questões discursivas, a média foi 32,8.
O mesmo ocorreu entre os Concluintes, que tiveram, em Formação Geral, média 41,0 nas questões objetivas e 33,3 nas questões discursivas.
Cabe destacar que, nas questões discursivas, o desempenho dos Concluintes foi muito pouco superior ao dos Ingressantes.
A média da nota máxima de Ingressantes, por exemplo, foi 95,0, superior à de Concluintes, 92,5, nesta parte do exame.
Os coeficientes de assimetria indicam assimetria positiva, de 0,1, para os Concluintes e para os Ingressantes.
Verifica -se a semelhança das distribuições das notas desses grupos, no Gráfico 3.9, apresentado a seguir.
Na sequência, os resultados verificados para cada uma das questões discursivas de Formação Geral serão apresentados, estabelecendo -se relações com os conteúdos englobados em cada uma delas.
Os comentários da banca corretora a respeito do observado nas respostas dos estudantes, suas impressões e conclusões serão apresentados junto à análise de cada questão.
Cumpre esclarecer que, tendo em vista que as questões discursivas de Formação Geral são padronizadas, ou seja, constam de todas as provas, os comentários da banca são os mesmos para todas as carreiras acadêmicas, sendo direcionados a todos os estudantes que participaram do ENADE/2010.
A seguir, serão analisados os desempenhos de Ingressantes e Concluintes da Área de Tecnologia em Agronegócio nas duas questões discursivas de Formação Geral do ENADE/2010. 3.3.1.1 Análise da Questão Discursiva 9 do Componente de Formação Geral Os dados obtidos a partir das respostas à questão 9 encontram -se na Tabela 3.13 e no Gráfico 3.10.
Nessa questão – de melhor desempenho dentre as duas de Formação Geral – os Ingressantes tiveram média 33,8, e os Concluintes, 35,3, resultando com diferença de 1,5 pontos entre os dois grupos.
Houve semelhança na distribuição de notas dos dois grupos, observando -se praticamente o mesmo percentual de Concluintes e Ingressantes que deixaram a questão em branco, cerca de 23,0 \%.
Quanto aos que obtiveram nota zero na questão ainda que não a tivessem deixado em branco, os percentuais foram de 4,0 \% para Ingressantes e de 2,8 \% para Concluintes.
Entre Ingressantes, a nota mediana foi 35,0 pontos, enquanto Concluintes obtiveram 40,0.
Estes últimos tiveram maior proporção de notas no intervalo de 50 a 60, com 13,6 \% dos estudantes nessa faixa.
A nota máxima, 100,0 pontos, foi alcançada tanto pelos Concluintes quanto pelos Ingressantes. 3.3.1.2 Comentários sobre as respostas à Questão Discursiva 9 No que tange mais precisamente à atualidade do tema proposto e ao desenvolvimento do conteúdo, deve ser ressaltado que se mostrou sintonizado com a realidade política brasileira, bem como com a necessária reflexão sobre o papel na sociedade do profissional formado pelas Instituições de Ensino Superior do país.
Dissertar sobre democracia e ética foi algo apropriado para a época em que a prova foi aplicada, pois coincidiu com o período das eleições, em que tais temas são mais discutidos na vida cotidiana.
Não se pode, entretanto, deixar de se levar em conta a tendência de os estudantes terem limitado a associação do tema “ ética ” ao assunto “ política ”, considerando -se que a ética abarca não apenas a política.
Tal associação atingiu, em alguma medida, o alcance mais preciso das respostas.
As respostas apresentadas refletiram motivação para escrever sobre o tema da questão.
Alguns estudantes chegaram a desenvolver textos eloquentes e engajados.
Em alguns casos, o conceito de democracia, por exemplo, foi bem desenvolvido ao longo do texto, simultaneamente às exemplificações pedidas.
Além disso, muitos estudantes associaram o exercício da ética ao seu futuro ofício, reconhecendo o seu papel como cidadão no desempenho de sua profissão.
Essa associação insinuou, também, o esforço de reconhecimento da noção de ética prática, ainda que de maneira codificada.
Um dos pontos positivos revelou -se no fato de, na maioria das respostas analisadas, haver uma preocupação em definir a sociedade democrática.
Isso facilitou a correção e elevou a média de notas.
Outro aspecto positivo foi a presença pertinente de exposições coerentes de comportamento ético de um futuro profissional.
A articulação dos itens “ a ” e “ b ” da questão produziu, portanto, todo tipo de reflexão acerca de uma postura ética na sociedade, inclusive, corporativa.
A comprovação de que o tema ética tem sido abordado nos cursos superiores foi outro destaque positivo detectado durante as correções.
Assim, além do caráter crítico de que muitas das respostas corretas se revestiram, os relatos constantes sobre desdobramentos da ética comprovaram a abordagem sobre o tema, no mínimo, a ética profissional, nas carreiras universitárias que participaram do ENADE.
A despeito de alguns estudantes se posicionarem contra o exame, seja anulando a questão ou escrevendo algum comentário desabonador sobre o processo, foi observado que significativa parte dos participantes tem consciência político-acadêmica e percebe por que está realizando aquele tipo de prova.
Desse modo, a maioria dos estudantes tinha algo a “ dizer “ sobre o tema da questão, isto é, valia -se da oportunidade para se posicionar a respeito do assunto trazido à baila.
Ainda do ponto de vista do conteúdo desenvolvido nas respostas, observou -se preocupação solidária e consciência cidadã dos estudantes.
Mais um destaque positivo crucial se refere à habilidade textual.
Os estudantes, por exemplo, expressaram atitudes éticas e não éticas e as contrapunham ao conceito de democracia.
Há casos em que os alunos definiam a corrupção política como um exemplo de atitude não ética, enfatizando a real missão desse cargo em uma sociedade democrática : “ atender aos anseios da população e zelar pelo bem comum ”, escreveram vários estudantes. de atitude não ética, enfatizando a real missão desse cargo em uma sociedade democrática : “ atender aos anseios da população e zelar pelo bem comum ”, escreveram vários estudantes.
No tocante à produção de texto dissertativo, à objetividade e ao caráter sintético, muitos estudantes corresponderam ao solicitado, elaborando pequenos textos dissertativos.
Nessas respostas, percebeu -se demonstração de domínio do assunto focado e utilização dos textos de referência como base para a elaboração.
O enunciado da questão 9 determinou que fosse elaborado um texto dissertativo que contivesse a )..., b )... e c )..., fornecendo um encaminhamento para as respostas que conduziu tanto a um resultado positivo quanto a um negativo, a saber : foi produtiva a resposta dos estudantes no sentido de que conseguiram, na maior parte das vezes, atender plenamente aos três quesitos apresentados na questão ; por outro lado, mais especificamente devido ao uso das três letras indicativas ( a, b e c ), houve tendência à elaboração de respostas divididas em itens, em vez de um texto eminentemente dissertativo.
Muitos estudantes apresentaram respostas nesse formato fragmentado para satisfazer a formalidade do roteiro sugerido no comando da questão.
A ideia-chave que deveria permear toda a resposta – democracia – se perdia em fragmentos desconexos entre si.
Se, por um lado, o roteiro do comando da questão facilitou o raciocínio de alguns dos estudantes, por outro, induziu grande parte deles a dois caminhos : i ) o das respostas curtas e objetivas, do tipo “ um profissional de saúde não ético é aquele que conta para as outras pessoas o que foi feito na consulta com o seu paciente ” e ii ) o outro caminho, que é do contraponto entre o conceito de sociedade democrática e a realidade brasileira.
Esse segundo caminho levou vários universitários a produzirem textos bem escritos, distantes, entretanto, do tema exigido.
No que se reporta às especificidades do desenvolvimento do conteúdo, a dificuldade ou falta de competência comunicativa revelada em grande parte das respostas denuncia a pouca exposição dos alunos à leitura de textos escritos de mais excelência.
Tal dificuldade mostrou -se mais evidente na impropriedade no uso de conectores e na incapacidade de elaboração de texto dissertativo, no tratamento de questões mais abstratas.
Muitos estudantes demonstraram não saber em que consistia a ética e confundiram exemplificação com definição : ao tentarem dar exemplos, via de regra, elaboravam uma definição incipiente do que se pedia.
Portanto, a dificuldade dos estudantes em argumentar e articular as ideias com coesão e coerência ainda são os grandes “ vilões ” desse processo avaliativo.
O excesso de fragmentos e textos prévios acerca do tema permitiu que a elaboração das respostas se resumisse a uma colagem ou paráfrase mal feita dos O excesso de fragmentos e textos prévios acerca do tema permitiu que a elaboração das respostas se resumisse a uma colagem ou paráfrase mal feita dos materiais disponíveis.
Em vez de subsidiar as respostas, esse material acabou por influenciar negativamente o desempenho dos estudantes.
Desse modo, verificou -se um nível de letramento ainda significativamente baixo, visto que a maioria apresentou uma visão bastante reducionista e simplificada sobre o tema em pauta.
De um modo geral, ao conceituarem ética e sociedade democrática, os estudantes demonstraram desconhecer o limite entre os dois conceitos, sem aludirem à interseção de ambos.
Foi expressiva a quantidade de redações nas quais questionavam a democracia pelo simples fato de estarem fazendo, a fortiori, aquela prova especificamente.
Em muitas das respostas, os estudantes se diziam incomodados pelo fato de estarem obrigados a realizar o exame, demonstrando sua insatisfação.
Por diversas vezes, os estudantes não foram capazes de sintetizar e de selecionar seu discurso o suficiente para atender ao que era pedido, o que ocasionou a falta de, pelo menos, um dos aspectos da resposta.
Foi constatada grande parte de respostas com problemas graves em termos textuais e gramaticais, em relação ao que se esperaria de pessoas que estão vinculadas ao Ensino ( Técnico ) Superior.
Alguns estudantes sequer conseguiam elaborar uma frase coerente e coesa, apresentando respostas que poderiam ser classificadas como de analfabetismo funcional.
De maneira ampla, manifestou -se da parte dos estudantes uma grande dificuldade no trato da língua escrita, dificuldade essa que se traduz, sobretudo, por um desrespeito bastante frequente às regras relativas : i ) à ortografia ( por exemplo : “ proficionais ”, “ infligir as leis ”, “ corrupitos ”, “ democrácia ”, “ opitar ”, “ manter o cizilo proficionau ” ), troca de “ sociedade democrática ” por “ sociedade demográfica ” e de “ ética ” por “ étnica ” ; ii ) à pontuação ; iii ) à concordância verbal ( “ os político deve (... ) eles pode ” ) ; iv ) ao modo de organização textual ( sem distinção entre o que seja dissertativo, argumentativo ou narrativo ).
Frisa -se que esse comentário não se restringe ao que se costuma chamar de “ língua padrão ”, mas ao que poderia ser chamado de “ excelência textual ”, que não se mede exclusiva ou enfaticamente pela correção gramatical, mas pela competência, objetividade, clareza e, em alguma medida, pela criatividade no ato de redigir.
Por fim, destaca -se um dos problemas dignos de nota na formulação da questão.
A transcrição do Dicionário Houaiss apresentava defeitos.
Por exemplo, a definição 5 de “ democracia ” indicava que esse é o tipo de “ governo que acata a vontade da maioria, embora respeitando os direitos e a livre expressão das minorias ”.
O emprego de um valor concessivo na oração adverbial, por oposição a um valor aditivo, contradiz as concepções efetivamente democráticas apresentadas de 1 a 4.
Repercussões desse valor concessivo, mal empregado, foram encontradas em incontáveis respostas, as quais diziam ser típico de um governo democrático, seguir a vontade da maioria, “ apesar ” de respeitar as minorias. um governo democrático, seguir a vontade da maioria, “ apesar ” de respeitar as minorias.
Outro desses problemas dignos de nota refere -se à indicação do dicionário quanto a “ governo ”, ao passo que o comando focava a ideia mais ampla de “ sociedade democrática ”. 3.3.1.3 Análise da Questão Discursiva 10 do Componente de Formação Geral O desempenho dos estudantes na questão 10 foi inferior ao obtido na questão de número 9, com nota média dos Ingressantes de 31,9, e de Concluintes, de 31,4.
A nota máxima também foi atingida em ambos os grupos, sendo que os percentuais de alunos que deixaram a questão em branco, tanto no grupo dos Ingressantes quanto no dos Concluintes, foram ligeiramente superiores àqueles apresentados na questão 9, 25,8 \% e 26,6 \%, respectivamente.
Entre as notas não nulas, houve maior percentual de alunos, considerando -se a distribuição de notas, no intervalo entre 40 e 50 pontos, tanto para Ingressantes ( 17,4 \% ) quanto para Concluintes ( 15,5 \% ).
Não se manifestou uma superioridade de Concluintes nas diferentes faixas de notas posteriores a esse intervalo, com equilíbrio no desempenho dos dois grupos, só ligeiramente superior no intervalo de notas entre 70 e 80 pontos.
Essas informações estão ilustradas na Tabela 3.14 e no Gráfico 3.11. 3.3.1.4 Comentários sobre as respostas à Questão Discursiva 10 Inicialmente, será apresentada uma panorâmica sobre as respostas, com destaque para os aspectos positivos observados.
Em relação ao desenvolvimento do tema e à vinculação do mesmo à vida cotidiana dos estudantes que participaram do exame, a visão positiva em relação à Educação de Jovens e Adultos ( EJA ) reforça a ideia de que a educação pode trazer crescimento não somente profissional, mas, sobretudo, para ampliar os horizontes da cidadania.
Foram frequentes as respostas que apontavam a educação como a melhor saída, senão a única, para o resgate da cidadania.
A atuação dos governos, em especial o governo federal, foi um ponto destacado nas respostas, as quais sinalizavam para a necessidade de as instâncias governamentais assumirem um papel cada vez mais ativo na consolidação da educação – pública, sobretudo – no país.
As políticas públicas focadas na educação foram, então, um item bastante frequente nas respostas, seja para enaltecer, criticar ou requisitar a ação governamental.
Outro ponto positivo destacado diz respeito à consciência ou ao reconhecimento da precariedade de condições de trabalho dos professores, em geral, no Brasil e, em particular, daqueles ocupados com turmas de jovens e adultos.
Nesse sentido, foram inúmeras as propostas apresentadas na segunda parte das respostas com relação à imperiosidade de melhoria salarial para os docentes, e à melhoria das condições gerais de trabalho desses profissionais.
A associação direta entre valorização do professor e sucesso escolar dos estudantes foi um ponto complementar a essa preocupação com os docentes.
Os estudantes demonstraram uma visão de vínculo no processo de ensino-aprendizagem, no qual alunos e professores são partes inalienáveis entre si.
A questão da repetência e da evasão escolar, referenciadas no suporte da questão, foram então, abordadas pelos alunos por meio do viés dessa valorização docente.
Ainda que os temas específicos da repetência e da evasão não fossem contemplados a contento na maior parte das respostas, ao menos, nessa associação com a qualificação e a valorização dos professores, algo foi trabalhado analiticamente pelos estudantes.
Num plano mais abrangente, houve reconhecimento de que, no processo ensino-aprendizagem, não há atores isolados : alunos, profissionais de educação, governos e a própria sociedade como corresponsáveis pelo sucesso ou fracasso de tal processo.
A atualidade do tema e a sua explícita vinculação com o cotidiano de muitos estudantes que realizavam o exame fizeram com que os respondentes demonstrassem um visível grau de envolvimento com o tema proposto.
Ou seja, foi possível detectar um empenho notório dos estudantes ao redigir suas respostas com seriedade e comprometimento relativo à delicada problemática dos milhões de estudantes trabalhadores do país.
Utilizar todo o espaço destinado à resposta, procurando aproveitar ao máximo a oportunidade de emitir um argumento ou uma opinião revelou -se como uma das evidências do aludido empenho.
A EJA também foi reconhecida como um dos caminhos concretos para o acesso à Universidade, no médio prazo.
As respostas se voltaram, muitas vezes, à condição de depoimentos pessoais sobre a conciliação de trabalho e estudo.
Nessa direção, ao assumirem o papel de depoentes, os estudantes expuseram suas histórias de vida, suas expectativas e percalços, suas esperanças e entraves cotidianos, seus anseios mais ambiciosos e problemas estruturais, que poderiam se resolver por meio da educação.
Nesse inciso, ficava clara a relação que se formava entre educação e mobilidade social, mediatizada por iniciativas governamentais como a EJA.
Uma sugestão plausível e válida para a solução dos problemas com a EJA no país reportou -se à necessidade de maior participação do público alvo na definição de estratégias para essa modalidade de ensino.
Nesse detalhe, era notória a associação das preocupações relativas à questão 9, que trata do problema da democracia e da cidadania, com o problema específico da educação escolar.
Houve um número significativo de sugestões ou propostas para se combater o problema da repetência e da evasão, destacando -se melhores salários para os Muitos estudantes abordaram, com propriedade, a associação entre um mercado de trabalho cada vez mais competitivo e a necessidade de qualificação profissional, inclusive com uma clara percepção de que os salários são tanto maiores quanto maior a quantidade de anos de estudo e/ou a qualidade da formação.
Nesse tópico, retorna -se ao papel crucial da EJA como viabilizadora dessa formação.
Em continuidade à análise, serão apresentados aspectos negativos que foram notados com frequência nas respostas à questão 10.
De um ponto de vista mais global, a dificuldade repousa no entendimento da questão, isto é, na decifração do comando.
Muitos fugiram completamente do assunto proposto, ainda que o enunciado da questão fosse farto em dados específicos ao tema.
Em vez de responderem diretamente quais eram os problemas enfrentados por aqueles que têm de trabalhar e estudar ao mesmo tempo, muitos estudantes optavam por falar das dificuldades que existem na educação brasileira, de forma estendida e pasteurizada.
Vários estudantes não conseguiam, em suas respostas, relacionar as dificuldades de conciliação da escola com o trabalho, a não ser, como já mencionado, nos casos de depoimentos pessoais, relatos de experiência, ou seja, traços de uma biografia.
Assim, eram respostas cujo teor era apenas o de registrar o abandono dos estudos por parte de vários jovens que têm de trabalhar.
O tom das respostas, nesses casos, era meramente de uma descrição e não de uma análise, muito menos de um elemento que compusesse uma argumentação.
A dificuldade de compreensão da questão se expressava nas respostas que, em lugar de mostrarem os problemas de conciliação entre escola e trabalho, prendiam -se à ideia de “ associação ” – termo que consta do enunciado da questão – entre escola e trabalho.
Muitos alunos relatavam as consequências positivas de uma boa educação no momento em que se busca bons postos de trabalho.
Tratava -se de uma fuga parcial ao comando da questão, evidenciando má compreensão do mesmo.
Uma dificuldade facilmente percebida nas respostas refere -se à tendência de os estudantes tratarem dos problemas relativos à escola e, muito separadamente, daqueles problemas relativos ao trabalho, reconhecendo uma espécie de sinergia que impediria de professores e aulas mais dinâmicas.
A lista dessas sugestões, contudo, era bem mais ampla e diversificada, incluindo ainda : a redução da jornada de trabalho dos jovens e adultos que participam da EJA ; melhoria de infraestrutura nas instituições responsáveis pela EJA ; incentivo financeiro aos estudantes trabalhadores por meio de um sistema de bolsas de estudo ; realização de campanhas motivacionais para o retorno aos estudos ; estímulo à participação das empresas, por meio da solicitação de recursos materiais, espaço físico para as salas, auxílio financeiro ; incremento dos cursos profissionalizantes no país ; e ampliação da oferta de ensino a distância em todo o território nacional.
Muitos estudantes abordaram, com propriedade, a associação entre um mercado de trabalho cada vez mais competitivo e a necessidade de qualificação profissional, inclusive com uma clara percepção de que os salários são tanto maiores quanto maior a quantidade de anos de estudo e/ou a qualidade da formação.
Nesse tópico, retorna -se ao papel crucial da EJA como viabilizadora dessa formação.
Em continuidade à análise, serão apresentados aspectos negativos que foram notados com frequência nas respostas à questão 10.
De um ponto de vista mais global, a dificuldade repousa no entendimento da questão, isto é, na decifração do comando.
Muitos fugiram completamente do assunto proposto, ainda que o enunciado da questão fosse farto em dados específicos ao tema.
Em vez de responderem diretamente quais eram os problemas enfrentados por aqueles que têm de trabalhar e estudar ao mesmo tempo, muitos estudantes optavam por falar das dificuldades que existem na educação brasileira, de forma estendida e pasteurizada.
Vários estudantes não conseguiam, em suas respostas, relacionar as dificuldades de conciliação da escola com o trabalho, a não ser, como já mencionado, nos casos de depoimentos pessoais, relatos de experiência, ou seja, traços de uma biografia.
Assim, eram respostas cujo teor era apenas o de registrar o abandono dos estudos por parte de vários jovens que têm de trabalhar.
O tom das respostas, nesses casos, era meramente de uma descrição e não de uma análise, muito menos de um elemento que compusesse uma argumentação.
A dificuldade de compreensão da questão se expressava nas respostas que, em lugar de mostrarem os problemas de conciliação entre escola e trabalho, prendiam -se à ideia de “ associação ” – termo que consta do enunciado da questão – entre escola e trabalho.
Muitos alunos relatavam as consequências positivas de uma boa educação no momento em que se busca bons postos de trabalho.
Tratava -se de uma fuga parcial ao comando da questão, evidenciando má compreensão do mesmo.
Uma dificuldade facilmente percebida nas respostas refere -se à tendência de os estudantes tratarem dos problemas relativos à escola e, muito separadamente, daqueles problemas relativos ao trabalho, reconhecendo uma espécie de sinergia que impediria de Isso significa que os estudantes distinguiam por mero isolamento, ou seja, eles não conseguiam distinguir e articular e, sobretudo, reconhecer as convocações mútuas entre os problemas da escola e do trabalho, no caso explícito dos estudantes trabalhadores.
Essa dificuldade de identificar, associar e analisar aspectos distintos, até mesmo concorrentes, antagônicos e complementares, denuncia a fragilidade de um pensamento simplificador que a qualidade das respostas revelou ser ainda hegemônico nas universidades brasileiras.
Notou -se a presença de preconceitos na análise dos fatores que levam ao abandono dos estudos.
Muitas respostas apontavam como causa desse abandono e da repetência a preguiça, a falta de perseverança e o cansaço dos estudantes.
Nesse tópico específico, fica evidente a responsabilização do indivíduo pelos percalços e fracassos de sua vida, descontextualizando -o de uma sociedade que, em seu movimento de conjunto, lhe impõe condicionantes.
A precariedade do uso da língua pátria de modo culto foi uma constante, conforme mencionado na análise da questão 9.
As impropriedades lexicais e a falta de elementos coesores que articulassem as partes constitutivas das respostas foram um incômodo constante.
Nesse sentido, agrava -se a situação da baixa qualidade textual em decorrência complementar da construção de textos meramente descritivos e não argumentativos, como é exigido no suporte da questão.
Isso demonstra a dificuldade dos estudantes universitários na compreensão de que uma argumentação requer, imperiosamente, justificações para os apontamentos apresentados.
Os textos das respostas eram, em sua maioria, opinativos, sem qualquer compromisso com elementos explicativos ou com a busca de justificativas. 3.3.2 Componente de Conhecimento Específico Na parte da prova relativa ao componente de Conhecimento Específico, observa -se que a diferença de desempenho foi mais acentuada entre Ingressantes e Concluintes.
A Tabela 3.15 contém dados elucidativos, na qual é possível observar que, no grupo dos Ingressantes, a nota média 14,6 é inferior à alcançada por Concluintes, 17,9.
Essas médias foram menores do que as obtidas no conjunto das questões objetivas, nas quais Ingressantes alcançaram a nota média 42,9, e Concluintes, 48,3 ( Tabela 3.8 ).
O Gráfico 3.12 apresenta a distribuição de notas nos diferentes intervalos.
Observa -se que Ingressantes e Concluintes atingiram notas mais elevadas na faixa entre zero e 10 ( 50,6 \% para os primeiros e 43,0 \% para os segundos ).
A faixa seguinte, entre 10 e 20 pontos, foi alcançada por 20,9 \% dos Ingressantes e 20,7 \% dos Concluintes.
Acima desses patamares, menciona -se o intervalo entre 20 e 30 pontos, com 14,4 \% e 14,7 \%, respectivamente.
A análise de cada uma destas questões será feita a seguir. ∗ ∗ A análise completa, com apresentação de trechos de respostas de estudantes, pode ser encontrada no documento “ Relatório da Correção das Questões Discursivas ”, elaborado pelos docentes que participaram do processo de correção das mesmas. 3.3.2.1 Análise da Questão Discursiva 38 do Componente de Conhecimento Específico Na questão 38, cujo desempenho aferido encontra -se descrito na Tabela 3.16, o grupo dos Ingressantes obteve a nota média 12,7, enquanto a média dos Concluintes foi 16,3.
Deixaram a questão em branco 26,8 \% dos Ingressantes e 23,6 \% dos Concluintes.
As concentrações de notas dos Concluintes foram superiores às das notas dos Ingressantes a partir do intervalo de 10 a 20 pontos.
A nota máxima foi 85,0 para Ingressantes e 80,0 para Concluintes. 3.3.2.2 Comentários sobre as respostas à Questão Discursiva 38 As respostas elaboradas pelos estudantes revelaram não haver compreensão do significado da expressão “ texto dissertativo ”.
A maioria apresentou textos com uma lista de itens, sem análise ou aprofundamento.
Destaca -se, ainda, que a maior parcela dos respondentes não demonstrou domínio do assunto, embora a questão possa ser classificada como de conhecimento geral.
A primeira parte da questão ofereceu mais oportunidades para que fossem dadas respostas com melhor elaboração.
Houve textos de estudantes que demonstraram imaturidade diante do tema.
Em casos de pontuação zero, os estudantes não responderam ao que foi pedido. 3.3.2.3 Análise da Questão Discursiva 39 do Componente de Conhecimento Específico A Tabela 3.17 contém as informações relativas à questão 39 do conjunto de questões do componente de Conhecimento Específico.
O desempenho dos dois grupos – Ingressantes e Concluintes – nessa questão, foi pouco superior ao desempenho na questão 38.
As médias foram 16,5 no grupo dos Ingressantes e 19,5 no grupo dos Concluintes.
A proporção de estudantes que deixaram a questão em branco foi maior do que a observada na questão 38, sendo esses percentuais 34,8 \% entre os Ingressantes e 31,8 \% entre os Concluintes.
A nota máxima ( 100,0 ) foi atingida pelos Ingressantes, ao passo que Concluintes atingiram 95,0.
Entre as notas não nulas, os intervalos de zero a 10 e de 20 a 30 pontos correspondem às notas de maior frequência nos dois grupos.
As distribuições das notas entre os estudantes constam do Gráfico 3.14. 3.3.2.4 Comentários sobre as respostas à Questão Discursiva 39 As respostas foram mais bem elaboradas que as da questão anterior.
Muitos estudantes, entretanto, confundiram “ Complexo agroindustrial ” com Agroindústria.
Em muitas das respostas, os estudantes se referiram a suas realidades locais, citando exemplos da região onde residem. 3.3.2.5 Análise da Questão Discursiva 40 do Componente de Conhecimento Específico A questão 40 da prova de Tecnologia em Agronegócio foi anulada pela comissão que representou a área junto ao INEP.
CONSIDERAÇÕES FINAIS A prova de Tecnologia em Agronegócio, segundo os professores corretores, caracterizou -se por apresentar temas atuais e do cotidiano, expostos pela imprensa brasileira.
Esperava -se, portanto, que os estudantes obtivessem notas maiores na média, uma vez que os assuntos eram de ampla divulgação.
Todas as questões eram realmente dissertativas, com comandos que favoreciam o processo de escrita.
No entanto, os estudantes, em sua maioria, não se expressaram com desenvoltura, tanto quanto à norma padrão da língua portuguesa quanto aos conteúdos específicos.
Em muitas situações, observou -se a falta de vocabulário pertinente da área, o que pode estar relacionado à falta de leitura.

% matched lemmas: ultrapassar
\end{textsample}
